% !TeX program = pdflatex
% !TeX encoding = UTF-8
\documentclass[a4paper,12pt]{article}
\usepackage[tikz]{relatorio}
\makeatletter\g@addto@macro{\UrlBreaks}{\UrlOrds}\makeatother  % permite quebrar URLs longas
\setstretch{1.1}
\renewcommand{\topfraction}{0.9}\renewcommand{\bottomfraction}{0.8}
\renewcommand{\textfraction}{0.1}\renewcommand{\floatpagefraction}{0.8}
\usepackage{enumitem}
\setlist{itemsep=2pt, topsep=4pt, parsep=0pt}

% ── DADOS DO RELATÓRIO ───────────────────────────────────────────────────────
\disciplina{Sistemas Paralelos e Distribuídos}
\professor{Ricardo Rios}

\titulo{Logística Distribuída de Soros Antivenenos em Regiões com Internet Instável}
\cabecalho{Logística distribuída de soros antivenenos}

\autor{Cristiano Peniche Ceccon}
\autor{João Victor Mattos Lopes}
\autor{Thiago Cordeiro de Melo}

\data{\today}

\descricao{%
  Entrega 1 -- Proposta do projeto. Documento de apoio para a discussão do
  escopo com o professor: problema, solução proposta e protótipo inicial.}

\begin{document}

\iniciarPreTextual
\fazcapa
\fazfolhaderosto

\iniciarTextual

% =============================================================================
\section{Proposta}
\label{sec:resumo}

Integrantes: Cristiano Peniche Ceccon, João Victor Mattos Lopes e Thiago
Cordeiro de Melo.

O soro antiveneno é o tratamento para acidentes com serpentes, escorpiões e
aranhas, e quanto antes for aplicado, melhor. Depois de seis horas as
complicações aumentam. O SUS distribui o soro de graça, mas na Amazônia um
paciente pode levar dias para chegar a um lugar que tenha o
soro~\cite{monteiro2020}. Para reduzir esse tempo, o Ministério da Saúde vem
levando o soro para postos menores e polos indígenas, mais perto das
comunidades~\cite{nt127,farias2025}.

Com isso o estoque deixa de ficar em poucos hospitais e passa a ficar
espalhado em muitos postos. Hoje não existe um
jeito confiável de saber onde tem soro, de qual tipo, de qual lote e até
quando ele vale. Os casos atendidos nesses postos também demoram a ser
notificados, o que atrasa a identificação de surtos e diminui a cota de soro
que a região recebe, porque a cota é calculada pelas notificações.

A conexão desses lugares melhorou muito com a internet via satélite de baixa
órbita. A Starlink já chegava a cerca de 90\% dos municípios da Amazônia Legal
em 2023~\cite{olhar2023}, e programas públicos levaram antenas a mais de mil
comunidades indígenas, quilombolas e ribeirinhas~\cite{olhar2024}. Na maioria
dos postos, porém, essa é a única conexão, e ela depende de energia (gerador
ou placa solar) e de um único fornecedor. Em 2024 a possível suspensão do
serviço no Brasil ameaçou desconectar centenas de comunidades da
Amazônia~\cite{teletime2024}, e em julho de 2025 uma falha de software derrubou
a rede no mundo todo por cerca de duas horas e meia~\cite{starlink2025}. Além
disso, a antena fica no posto, enquanto o acidente acontece no rio, na roça ou
na aldeia.

Por isso o registro de uma dispensação ou de um caso não deve depender de um
servidor central acessível. Cada posto é um nó do sistema: quando a conexão
cai, ele continua funcionando, e quando ela volta seus dados são juntados com
os dos outros sem perder nem duplicar nada. Com internet, essa sincronização
leva segundos; sem ela, cada informação mostra há quanto tempo foi
atualizada.

Para isso, cada posto vai guardar uma cópia local dos dados e sincronizar com
os vizinhos sempre que tiver conexão. O foco do projeto, no entanto, são as
cinco funções abaixo, que seriam úteis mesmo com internet perfeita. A cópia
local serve para que elas continuem disponíveis quando a conexão falhar:
\begin{enumerate}
  \item controle de estoque por posto, tipo de soro e lote, usando primeiro o
        lote que vence antes, respeitando o espaço refrigerado de cada posto e
        mostrando há quanto tempo cada informação foi atualizada;
  \item notificação de casos, que fica gravada no posto e vai para a central
        nas sincronizações;
  \item alerta de surto por região: o polo e a central comparam os casos do
        dia com a média dos dias anteriores (método C2 do CDC~\cite{hutwagner2003})
        e, se houver aumento fora do normal, avisam os postos daquela região e
        reforçam o estoque deles;
  \item indicação do posto mais próximo que tem o soro necessário, levando em
        conta barco, voadeira, estrada e avião;
  \item identificação dos postos críticos, isto é, os que deixariam
        comunidades a mais de seis horas do tratamento se ficassem sem soro.
        Cada posto é um cálculo independente, então isso roda em paralelo.
\end{enumerate}

Ao final esperamos ter o sistema rodando com cada nó em um contêiner separado e
avaliá-lo pelo tempo até todos os nós terem os mesmos dados, pelo atraso das
notificações e dos alertas, pelo volume de dados trocados, pelas decisões
tomadas com informação desatualizada e pelo ganho do paralelismo. O sistema não
faz diagnóstico nem indica tratamento.

% =============================================================================
\section{Contexto}
\label{sec:contexto}

Entre 2007 e 2019 foram notificados quase 1,9 milhão de acidentes por animais
peçonhentos no Brasil (Tabela~\ref{tab:sinan}).

\begin{table}[H]
  \centering
  \caption{Acidentes por animais peçonhentos notificados no SINAN, 2007--2019~\cite{completude2023}.}
  \label{tab:sinan}
  \begin{tabular}{lrr}
    \toprule
    Tipo de acidente & Casos & Proporção \\
    \midrule
    Escorpiões & 1.111.300 & 59,4\% \\
    Serpentes  &   386.938 & 20,7\% \\
    Aranhas    &   373.224 & 19,9\% \\
    \midrule
    Total & 1.871.462 & 100\% \\
    \bottomrule
  \end{tabular}
\end{table}

No país como um todo o escorpião é o mais comum, mas no Norte o problema maior
são as serpentes. A região concentrou 31,6\% dos acidentes com serpentes do
país, e no Amazonas cerca de 90\% deles são causados pela jararaca-da-Amazônia
(\emph{Bothrops atrox})~\cite{diversitas2022}.

O projeto SAVING, da Fundação de Medicina Tropical (ligada à UEA) com o
Instituto Butantan e o Ministério da Saúde, levou o soro a 14 polos de saúde
indígena no Amazonas~\cite{farias2025}. Segundo o Butantan, a proporção de
pacientes tratados nas primeiras seis horas subiu de 40,8\% para
75,8\%~\cite{butantan40mi}. Em maio de 2026 a Nota Técnica Conjunta
nº~127/2026~\cite{nt127} padronizou os Pontos Estratégicos de Antivenenos
(PEA), que já eram cerca de 2.319 unidades no país.

A mesma nota técnica pede que as unidades com estoque acompanhem a validade e
informem consumo, perdas e necessidade de reposição, e que a lista de PEA seja
mantida atualizada para o SAMU e os centros de toxicologia~\cite{nt127}. Saber
onde tem soro já é uma exigência oficial, o difícil é manter isso confiável em
centenas de pontos remotos. Existe ainda um ciclo vicioso: onde o acesso é
pior, notifica-se menos~\cite{completude2023}, e quem notifica menos recebe
menos soro.

Outro ponto é que soros e vacinas do Programa Nacional de Imunizações ficaram
de fora do identificador único por embalagem criado pela Anvisa~\cite{rdc157},
então o soro só é acompanhado por lote. A distribuição também não segue um
padrão nacional: o Ministério manda cotas para os estados e cada estado
combina seus próprios fluxos~\cite{nt127}, sem nada que obrigue um centro maior
a abastecer os menores. Além disso, só pode receber soro a unidade que tem
espaço livre na rede de frio.

% =============================================================================
\section{Solução proposta}
\label{sec:solucao}

Quando a rede cai, um sistema distribuído tem que escolher entre recusar
operações, para que todos os nós tenham sempre os mesmos dados, ou aceitar
operações e deixar os dados diferentes por um tempo (teorema
CAP~\cite{gilbert2002}). Um posto não pode deixar de registrar o soro que
aplicou porque está sem internet, então ficamos com a segunda opção.

Para juntar os dados depois sem conflito, vamos usar CRDTs (tipos de dados
replicados livres de conflito)~\cite{shapiro2011}. São estruturas pensadas
para que a junção de duas cópias dê sempre o mesmo resultado, não importa a
ordem em que as atualizações chegam nem se chegam repetidas. No estoque de um
lote, por exemplo, em vez de guardar só o saldo, cada cópia guarda quanto
entrou e quanto saiu por cada dispositivo. Ao juntar duas cópias, fica o maior
valor de cada dispositivo, e assim uma saída nunca é contada duas vezes. Casos
e alertas formam um conjunto que só cresce, com um identificador único para
cada item, então receber o mesmo caso duas vezes não muda nada.

De tempos em tempos cada nó escolhe um vizinho com quem tem conexão e os dois
trocam e juntam seus dados (anti-entropia ou \emph{gossip}~\cite{demers1987}).
Não existe caminho fixo: os dados de um posto isolado chegam à central por
qualquer vizinho conectado. A Figura~\ref{fig:arquitetura} mostra a topologia
usada no protótipo.

\subsection{Limitações}

Enquanto um posto está isolado, os outros continuam vendo o estoque dele da
última sincronização. Isso não tem como ser eliminado, por isso o sistema
mostra a idade de cada informação, e a ideia é medir esse efeito e
reduzi-lo. Também pode acontecer de dois aparelhos isolados no mesmo posto
registrarem a saída do mesmo frasco, deixando o saldo negativo. O sistema
detecta esse caso, mas para evitá-lo seria preciso definir antes quanto cada
aparelho pode retirar~\cite{balegas2015}. Por fim, os dados compartilhados
entre os nós não identificam o paciente (LGPD), e o registro clínico continua
nos sistemas oficiais.

% =============================================================================
\section{Trabalhos relacionados}
\label{sec:relacionados}

A Tabela~\ref{tab:relacionados} compara a proposta com os sistemas e estudos
mais próximos. A parte teórica se apoia em CRDTs~\cite{shapiro2011}, em
algoritmos epidêmicos de sincronização~\cite{demers1987} e em redes
tolerantes a atraso, nas quais até um barco pode levar dados de uma comunidade
para outra~\cite{pentland2004}.

\begin{table}[htbp]
  \centering
  \caption{Comparação com sistemas e estudos existentes.}
  \label{tab:relacionados}
  \small
  \begin{tabularx}{\textwidth}{>{\raggedright\arraybackslash}p{2.7cm}XX}
    \toprule
    Trabalho & O que faz & Diferença para esta proposta \\
    \midrule
    SIES e SINAN (Ministério da Saúde)~\cite{nt127} & Registros oficiais de
      insumos e de notificações. & São centralizados e dependem de conexão ou
      de digitação posterior. A proposta funciona sem internet e pode
      alimentá-los. \\
    mSupply~\cite{msupply} & Controle de estoque de medicamentos que funciona
      sem internet. & Sincroniza apenas com um servidor central; não trata
      surtos nem rotas. \\
    DHIS2~\cite{dhis2} & Coleta de dados de saúde com aplicativo que funciona
      sem internet. & Voltado a indicadores, não a estoque nem a decisões
      logísticas. \\
    Estudo da Fiocruz~\cite{fiocruz2024} & Estima o tempo de acesso a serviços
      com soro por tipo de transporte. & É uma análise pontual; a proposta usa
      o estoque atual de cada posto. \\
    \bottomrule
  \end{tabularx}
\end{table}

% =============================================================================
\section{Protótipo e primeiros resultados}
\label{sec:prototipo}

Implementamos um protótipo em Python, sem bibliotecas externas (pasta
\texttt{prototipo/}, executado com \texttt{python3 demo.py}). Ele cria oito
nós que se comunicam por TCP e simula 30 dias: acidentes, aplicação de soro,
remessas, conexões que caem e voltam, nove dias sem comunicação em Atalaia do
Norte e um surto no Alto Solimões entre os dias 14 e 19. É um cenário
propositalmente pessimista, para testar o sistema nas piores condições. Estoques, tempos e
taxas de acidentes são ilustrativos. A Tabela~\ref{tab:demo} resume uma
execução.

\begin{table}[htbp]
  \centering
  \caption{Resultados de uma execução do protótipo.}
  \label{tab:demo}
  \small
  \begin{tabularx}{\textwidth}{Xr}
    \toprule
    Medida & Valor \\
    \midrule
    Casos simulados / notificações perdidas & 111 / 0 \\
    Atraso típico (mediana) da notificação até a central & 7,1 h \\
    Maior atraso de notificação (Atalaia do Norte, isolada) & 180 h \\
    Atraso do alerta de surto até os postos conectados da região & 4,0 h e 13,6 h \\
    Atraso do alerta de surto até Atalaia do Norte (isolada) & 143,5 h \\
    Idade da informação de Atalaia vista pela central no dia 18 & 6,2 dias \\
    Tempo para todos os nós terem dados idênticos após a rede voltar & 8 ms \\
    Remessas acima do espaço refrigerado do destino & 0 de 27 \\
    Pacientes encaminhados a um posto que já estava sem soro & 1 de 9 \\
    \bottomrule
  \end{tabularx}
\end{table}

\begin{figure}[htbp]
  \centering
  \begin{tikzpicture}[
      scale=0.8, every node/.style={transform shape},
      font=\sffamily\footnotesize,
      no/.style={draw, rounded corners=2pt, minimum width=2.25cm, minimum height=0.8cm,
                 align=center, fill=#1},
      no/.default=white,
      bom/.style={thick},
      ruim/.style={thick, dashed}]
    \node[no=blue!12]  (man) at (0,0)      {Central\\Manaus};
    \node[no=green!12] (tab) at (-3.4,-1.9) {Polo\\Tabatinga};
    \node[no=green!12] (tef) at (3.9,-1.9)  {Polo\\Tefé};
    \node[no] (bc)  at (-5.6,-3.9) {Posto\\Benjamin Constant};
    \node[no] (ata) at (-3.4,-5.4) {Posto\\Atalaia do Norte};
    \node[no] (spo) at (-1.3,-3.9) {Posto\\S. Paulo de Olivença};
    \node[no] (alv) at (2.5,-3.9)  {Posto\\Alvarães};
    \node[no] (uar) at (5.3,-3.9)  {Posto\\Uarini};
    \draw[bom] (man) -- (tab);
    \draw[bom] (man) -- (tef);
    \draw[ruim] (tab) -- (tef);
    \draw[bom] (tab) -- (bc);
    \draw[ruim] (tab) -- (ata);
    \draw[ruim] (bc) -- (ata);
    \draw[ruim] (tab) -- (spo);
    \draw[bom] (tef) -- (alv);
    \draw[ruim] (tef) -- (uar);
    \draw[ruim] (alv) -- (uar);
    \node[draw, dashed, rounded corners, fit=(tab)(bc)(ata)(spo), inner sep=6pt,
          label={[font=\sffamily\scriptsize]below:Região Alto Solimões}] {};
    \node[draw, dashed, rounded corners, fit=(tef)(alv)(uar), inner sep=6pt,
          label={[font=\sffamily\scriptsize]below:Região Médio Solimões}] {};
    \begin{scope}[shift={(2.0,-5.6)}]
      \draw[bom] (0,0) -- (0.7,0) node[right] {conexão quase sempre disponível};
      \draw[ruim] (0,-0.45) -- (0.7,-0.45) node[right] {conexão instável};
    \end{scope}
  \end{tikzpicture}
  \caption{Topologia do protótipo. Cada linha liga dois nós que trocam dados
  diretamente quando há conexão. A central abastece os polos, e os polos
  abastecem os postos da sua região. Os municípios são reais; conexões e
  tempos são ilustrativos.}
  \label{fig:arquitetura}
\end{figure}

Nenhuma notificação se perdeu e todos os nós ficaram com os mesmos dados logo
que a rede voltou. Na primeira versão do protótipo, como o polo não via o
consumo de Atalaia, ele repetia remessas, e o posto chegou a quase quatro vezes
o estoque inicial. Corrigimos isso fazendo a remessa contar como estoque do
destino até a confirmação de recebimento chegar e limitando cada remessa ao
espaço refrigerado livre. Cada nó também passou a mostrar há quanto tempo
recebeu dados de cada posto. Ainda há decisões tomadas com informação velha,
como o paciente encaminhado para um lugar que já estava sem soro, e medir e
reduzir esses casos é o foco das próximas entregas.

Na parte paralela, a busca de postos críticos em um mapa sintético com 22.500
localidades e 200 postos levou 13,0~s com um processo e 3,9~s com 16 (3,3
vezes mais rápido). O ganho foi limitado pelo processador do notebook usado
(Intel i7-12700H), que reduz a frequência sob carga.

% =============================================================================
\section{Próximos passos}
\label{sec:discussao}

Os próximos passos são rodar cada nó em um contêiner próprio, com
quedas de rede reais e dados gravados em disco; enviar só o que mudou em cada
sincronização; e montar a avaliação com as medidas da
Seção~\ref{sec:resumo}. A avaliação também vai incluir um cenário mais
próximo da conexão via satélite, com alta disponibilidade e falhas raras
(falta de energia, chuva forte, pane do provedor), e uma comparação com um
sistema centralizado, contando quantas operações ficariam bloqueadas em cada
queda.

% =============================================================================
\footnotesize
\begin{thebibliography}{99}
\addcontentsline{toc}{section}{Referências}
\setlength{\itemsep}{0pt}

\bibitem{monteiro2020} MONTEIRO, W. M. et al. Providing antivenom treatment
  access to all Brazilian Amazon indigenous areas: ``Every life has equal
  value''. \textit{Toxins}, 2020.

\bibitem{nt127} BRASIL. Ministério da Saúde. \textit{Nota Técnica Conjunta
  nº 127/2026-SVSA/SAPS/SAES/SESAI/MS}: escolha dos Pontos Estratégicos de
  Antivenenos (PEA). Brasília, maio 2026.

\bibitem{farias2025} SEABRA DE FARIAS, A. et al. Decentralization of
  snakebite antivenom treatment to indigenous community health centers in the
  Brazilian Amazon (the SAVING Program). \textit{PLoS Neglected Tropical
  Diseases}, 2025.

\bibitem{diversitas2022} Epidemiologia de acidentes ofídicos no estado do
  Amazonas entre 2010--2020. \textit{Diversitas Journal}.

\bibitem{completude2023} Completude das notificações dos acidentes por
  animais peçonhentos no SINAN: estudo descritivo, Brasil, 2007--2019.
  \textit{Epidemiologia e Serviços de Saúde}, 2023.

\bibitem{butantan40mi} INSTITUTO BUTANTAN. Envio de soro antiofídico para
  regiões remotas da Amazônia pouparia R\$ 40 milhões (notícia).

\bibitem{gilbert2002} GILBERT, S.; LYNCH, N. Brewer's conjecture and the
  feasibility of consistent, available, partition-tolerant web services.
  \textit{SIGACT News}, 2002.

\bibitem{shapiro2011} SHAPIRO, M. et al. Conflict-free replicated data
  types. \textit{SSS}, LNCS 6976, 2011.

\bibitem{demers1987} DEMERS, A. et al. Epidemic algorithms for replicated
  database maintenance. In: \textit{ACM PODC}, 1987.

\bibitem{hutwagner2003} HUTWAGNER, L. et al. The bioterrorism preparedness
  and response Early Aberration Reporting System (EARS). \textit{J. Urban Health}, 2003.

\bibitem{balegas2015} BALEGAS, V. et al. Extending eventually consistent
  cloud databases for enforcing numeric invariants. \textit{IEEE SRDS}, 2015.

\bibitem{rdc157} BRASIL. Anvisa. \textit{Resolução RDC nº 157, de 11 de maio
  de 2017}: implantação do Sistema Nacional de Controle de Medicamentos.

\bibitem{olhar2023} Internet da Starlink já está em 90\% dos municípios da
  Amazônia Legal. \textit{Olhar Digital}, 23 out. 2023.

\bibitem{olhar2024} Indígenas na Amazônia recebem internet da Starlink.
  \textit{Olhar Digital}, 6 jun. 2024.

\bibitem{teletime2024} Possível suspensão da Starlink pode desconectar
  centenas de comunidades tradicionais da Amazônia. \textit{Teletime}, 2 set. 2024.

\bibitem{starlink2025} Starlink software failure causes global outage.
  \textit{Broadband Breakfast}, jul. 2025.

\bibitem{msupply} mSUPPLY FOUNDATION. \textit{Open mSupply: synchronisation
  (offline first)}. 2024.

\bibitem{dhis2} DHIS2. \textit{DHIS2 Documentation}. 2024.

\bibitem{fiocruz2024} Estudo analisa acesso a soro antiofídico em regiões
  com alto índice de acidentes com cobras. EBC.

\bibitem{pentland2004} PENTLAND, A.; FLETCHER, R.; HASSON, A. DakNet:
  rethinking connectivity in developing nations. \textit{IEEE Computer},
  v.~37, n.~1, 2004.

\end{thebibliography}

\end{document}
